% Generated draft: requires mathematical and source review.
\documentclass[11pt]{article}
\usepackage[margin=25mm]{geometry}
\usepackage{fontspec}
\setmainfont{Latin Modern Roman}
\usepackage{amsmath,amssymb,mathtools,unicode-math}
\setmathfont{Latin Modern Math}
\usepackage{enumitem,needspace,xurl,hyperref}
\setlength{\parindent}{0pt}
\setlength{\parskip}{5pt}
\setlength{\emergencystretch}{3em}
\newcommand{\N}{\mathbb{N}}
\newcommand{\Z}{\mathbb{Z}}
\newcommand{\Q}{\mathbb{Q}}
\newcommand{\R}{\mathbb{R}}
\newcommand{\C}{\mathbb{C}}
\newcommand{\F}{\mathbb{F}}
\newcommand{\A}{\mathbb{A}}
\newcommand{\PP}{\mathbb{P}}
\newcommand{\E}{\mathbb{E}}
\newcommand{\eps}{\varepsilon}
\newcommand{\dd}{\,\mathrm{d}}
\begin{document}
\section*{UnsolvedMath ID: 30000794}
\textbf{Source problem code:} catalogue entry
\subsection{Definitions and mathematical statement}
Let \(K\) be the function field of the curve over \(\F_q\) in the source, and let \(n_K(D)\) count the intended \(S_5\)-extensions with squarefree discriminant \(D\). The two predictions are Poisson mean \(\mu_K=1\) in this function-field setting and, for the analogous count \(n_\Q(D)\),
\[
 \mu_\Q=\frac{13}{120}\zeta(2)^{-1},
 \qquad \zeta(2)=\sum_{j\geq1}j^{-2}.
\]
The numerical factors and the distinction between the two means are retained exactly.

A deterministic count has a distribution only after specifying how \(D\) varies. If \(\mathcal D_B\) is the intended finite sampling set of squarefree discriminant data and \(n(D)\) the corresponding count, the Poisson assertion would mean
\[
 \forall j\in\Z_{\geq0},\qquad
 \frac{|\{D\in\mathcal D_B:n(D)=j\}|}{|\mathcal D_B|}
 \longrightarrow e^{-\mu}\frac{\mu^j}{j!}
\]
in the source's limiting regime, using the relevant \(\mu\) above. The class \(\mathcal D_B\), weighting and limit are absent. Sampling all admissible squarefree data, including those with count zero, is different from conditioning on a positive count. For function fields the degree/divisor ordering, fixed \(q\) versus growing \(q\), curve and local restrictions need specification. Over \(\Q\), sign, signature, discriminant ordering and local conditions are missing. No one regime is chosen silently.

The term \(S_5\)-extension is also unspecified. A degree-five separable extension with normal closure group \(S_5\), and a Galois extension with group \(S_5\) and degree \(120\), are different counted objects with different discriminants. The squarefree hypothesis does not resolve this convention by itself. One must state which extension discriminant, ideal/divisor or integer, is meant and whether objects are counted up to base-field isomorphism or with additional markings. In a divisor/ideal convention squarefree means all nonzero prime valuations of the effective discriminant are one.

The record additionally asks for experimental testing. In a specified model this entails enumerating the relevant counts and comparing empirical frequencies, including zeros, with the stated Poisson probabilities; it does not provide data, an enumeration algorithm or an error guarantee. This request is retained without reporting an experiment. The literature assessment checked 2026-08-21 refers to random number fields, whereas the cleaned formulation varies discriminant data over a stated base field; it gives no common probability model. Its report that related nonabelian Cohen--Lenstra work does not establish the specific law, the open label and source-verification note remain dated metadata.
\subsection{Short English statement}
Test whether the specified extension counts over varying squarefree discriminants have Poisson laws, with mean one over the stated function field and mean \((13/120)\zeta(2)^{-1}\) over \(\Q\). The extension convention, sampling of discriminants and limiting regime are essential missing data. Experimental testing remains an explicit requested task, with no experiment claimed here.
\subsection{Sources}
\begin{itemize}
\item Local UnsolvedMath registry, permanent ID 30000794, source code catalogue entry, ``Poisson Statistics for S5 Number-Field Extensions,'' including both predicted means, function-field formulation and experimental request.
\item \emph{Explicit Methods in Number Theory (2008)}, Oberwolfach Reports, pp.\ 1957--2034, DOI 10.4171/owr/2007/34: \url{https://doi.org/10.4171/owr/2007/34}. The citation uses 2008, while the DOI and proposed-year metadata use 2007.
\item Supplied reference ``current survey of random arithmetic groups'': \url{https://arxiv.org/abs/2301.09687}.
\item Literature triage checked 2026-08-21; record updated 2026-09-16. No external source has been visited or independently checked here.
\end{itemize}
\end{document}
